% !TeX root = ../papers/nc-p3.tex
\section{任意のトーラス作用に対する四文字のEuler行列式}
\label{sec:torus-euler-characters}
四元群印・点配置・領域性を仮定しない。標準Koszul左右AS四窓を保つトーラスの
作用を各 $V_i$ に線形に持ち上げる。$A_i$ を四矢、$B_i$ を六関係の指標とし、
$h$ を最上四テンソルの一次元指標とする。以下、$h$ は指標環の単項式、
$A^*$ は指数を逆にした指標とする。始点を法4で並べる。
\begin{theorem}\label{thm:torus-euler-characters}
Koszul分解の行列 $Q(x)$ の第 $i$ 行には、巡回順に
\[
 1+x^4h,\qquad -x A_i,\qquad x^2 B_i,\qquad -x^3 h A_{i+3}^*
\]
が入る。四つのトーラス文字 $m_1,\dots,m_4$ が重複を込めて一意に存在し、
\[
 \det Q(x)=\prod_{j=1}^4(1-x^4m_j),\qquad
 \prod_j m_j=h^4,\qquad
 \prod_{i=0}^3\det(V_i)=h^4 .
\]
最後の等式は射影安定群の線形逆像の恒等成分上での指標の等式でもある。
とくに四位置で行列式1の連続安定群は最上テンソルそのものを固定する。
\end{theorem}
\begin{proof}
まず一つの乗法群で文字変数を $t$ とする。$Q^{-1}$ は始点・終点の組ごとの
経路の指標Hilbert行列である。$|t|=1$ なら各係数の絶対値は
$\binom{n+3}{3}$ 以下であり、全成分が $|x|<1$ で正則である。
従って $Q(x)$ の行列式はこの円板で消えない。置換展開の各項の経路長の和は4の倍数なので
$y=x^4$ とおけば行列式 $D(y,t)$ は $y$ の四次式で、定数項1、最高次係数 $h^4$ である。
$|t|=1$ では両端係数の絶対値は1なので四つの逆根 $\lambda_j(t)$ は全て単位円にある。
$t=1$ では $Q=(I-xS)^4$（$S$ は巡回置換）より $D(y,1)=(1-y)^4$。

ここで四次多項式の逆根が実際に単項式になることを示す。
べき和 $p_n(t)=\sum_j\lambda_j(t)^n$ はNewton恒等式で整数Laurent多項式であり、
単位円上の絶対値が4以下なので各Fourier係数の絶対値も4以下である。
素数 $\ell=11$ を使う。Frobenius合同
$p_{n\ell}(t)\equiv p_n(t)^\ell\equiv p_n(t^\ell)\bmod\ell$ と係数の上界から、
$p_{n\ell}(t)=p_n(t^\ell)$ が全 $n$ で成立する。従って $t^\ell$ の逆根の多重集合は
$t$ の逆根の $\ell$ 乗の多重集合である。$t=\exp u$ とおき、$u=0$ でPuiseux級数の
枝を全て取る。全逆根の定数項は1なので形式対数 $g_j(u)$ が取れる。
多重集合 $\{g_j(\ell u)\}$ と $\{\ell g_j(u)\}$ の一致を与える置換を一つ固定し、
その位数 $N$ 回反復すると各枝で $g_j(\ell^N u)=\ell^N g_j(u)$ となる
（分数冪には正の実数の同じ根を使う）。係数比較により $g_j(u)=a_j u$。
$p_1(e^u)$ は有限個の $e^{nu}$, $n\in\mathbb Z$ の線形結合であり、
異なる指数の指数関数は原点の導関数のVandermonde行列により一次独立である。
従って全 $a_j$ は整数で、$p_1(t)$ の係数は四文字の多重度である。
トーラスが高次元なら $p_1$ の台およびその四重までの和と $D$ の全係数の台を
分離する一般の整数一径数を選ぶ。この方向に制限した係数比較は元の指標環にも持ち上がる。

最後に、$p_1=-[y]D$ を展開すると
\[
 p_1=\prod_i A_i-\sum_i B_i A_{i+2}A_{i+3}
       +h\sum_i A_i A_i^*+B_0B_2+B_1B_3-4h .
\]
任意の一径数で $h=t^H$ とし、$M=\sum_i A_i'(1)$ と置く。
$B_{i+2}=hB_i^*$ より $\sum_i B_i'(1)=12H$。
上式を微分すると $p_1'(1)=16M-60H$ である。一方、四文字の積は最高次係数
$h^4$ だから $p_1'(1)=4H$。従って $M=4H$ で、これが所望の行列式指標である。
連結線形群の文字は極大トーラスで決まることから最後の主張を得る。
\end{proof}

\begin{corollary}\label{cor:two-weight-profiles}
各位置の非自明な重みが二つまでで、スカラーを引いて重み $0,1$ となる一径数作用を考える。
スカラーだけの位置の重みは0とし、少なくとも一位置が非自明とする。
$k_i$ を重み1の矢次元、$H$ を最上の重みとする。可能な型は巡回を含め下表に限られる。
$B_i$ は重み $0,1,2$ の多重度順で、$B_{i+2}(t)=t^HB_i(t^{-1})$ である。
\[
\begin{array}{c|c|cc|c}
 (k_i)&H&B_0&B_1&\{\deg m_j\}\\ \hline
 1111&1&330&330&0004\\
 2222&2&141&141&0044\\
 3333&3&033&033&0444\\
 0121&1&420&240&0022\\
 0323&2&240&042&1133\\
 1133&2&321&060&2222\\
 1133&2&420&141&2222
\end{array}
\]
これは必要条件である。特に各位置が $2+2$ なら各辺に低位二文字の関係と
高位二文字の関係がちょうど一つずつ存在し、混合関係はちょうど四つである。
この場合の二重Hilbert級数は全始点で $(1-u)^{-2}(1-v)^{-2}$ である。
\end{corollary}
\begin{proof}
一文字全階数と重み0の存在から $1\le H\le3$。各辺の $B_i$ の重み別次元は
その辺の重み空間と、対辺の相補重み空間の両次元以下である。$p_1$ の上式から
四指数は全て0から4の間にあり、その和は $4H$。
\path{computations/verify_torus_euler_profiles.py} はこの範囲の整数多重度を尽くす。
行・列の経路長の差を取り、行列式には継ぎ目をまたぐ成分だけ周期変数 $y$ を付ければよい。
整数点 $(t,y)=(11,7)$ で全因子候補と比較して不一致を落とし、残った各組では
多項式の両変数の次数上界を超える格子で係数の一致も確かめる。
表のほかには $0103,H=1,B_0=B_1=510$ のみである。この型では位置1の高位文字を $v$ として、
両隣が重み0のみなので $a\otimes v$ と $v\otimes b$ が続けて関係に入り、
定理\ref{thm:no-consecutive-pure-relations}に反する。
$2222$ のHilbert級数は位置によらない上記の矢・関係の多重度をKoszul逆公式に入れると得られる。
\end{proof}
